\documentclass[conference]{IEEEtran}
\IEEEoverridecommandlockouts
% The preceding line is only needed to identify funding in the first footnote. If that is unneeded, please comment it out.
%Template version as of 6/27/2024

\usepackage{cite}
\usepackage{amsmath,amssymb,amsfonts}
\usepackage{algorithmic}
\usepackage{graphicx}
\usepackage{textcomp}
\usepackage{xcolor}

\usepackage{tikz}
\usetikzlibrary{shapes,arrows}

\usepackage{listings}
\usepackage{color}

\definecolor{dkgreen}{rgb}{0,0.6,0}
\definecolor{gray}{rgb}{0.5,0.5,0.5}
\definecolor{mauve}{rgb}{0.58,0,0.82}

\lstset{frame=tb,
  language=Python,
  aboveskip=3mm,
  belowskip=3mm,
  showstringspaces=false,
  columns=flexible,
  basicstyle={\small\ttfamily},
  numbers=none,
  numberstyle=\tiny\color{gray},
  keywordstyle=\color{blue},
  commentstyle=\color{dkgreen},
  stringstyle=\color{mauve},
  breaklines=true,
  breakatwhitespace=true,
  tabsize=3
}
\def\BibTeX{{\rm B\kern-.05em{\sc i\kern-.025em b}\kern-.08em
    T\kern-.1667em\lower.7ex\hbox{E}\kern-.125emX}}
\begin{document}



\tikzstyle{block} = [draw, fill=white, rectangle, 
    minimum height=3em, minimum width=6em]
\tikzstyle{sum} = [draw, fill=white, circle, node distance=1cm]
\tikzstyle{input} = [coordinate]
\tikzstyle{output} = [coordinate]
\tikzstyle{pinstyle} = [pin edge={to-,thin,black}]

\title{EEE 302 Project\\
    \author{\IEEEauthorblockN{1\textsuperscript{st} Tolga Uçar}
    \IEEEauthorblockA{\textit{Izmir University of Economics} \\
    Izmir, Türkiye \\
    freetolga@protonmail.com}
    }
}
\maketitle

\begin{abstract}
    This project is a techical overview on the inner-workings of frequency modulation(FM) and implementation in discrete-time using Python ,scipy and matplotlib
\end{abstract}

\begin{IEEEkeywords}
    signalprocessing, python, frequencymodulation, beginner
\end{IEEEkeywords}

\section{Introduction}

This paper will go over how frequency modulation works and examples of implementing it in the discrete world

\subsection{Definition of FM}

FM is defined as varying (modulating) the frequency of a sinusoidal carrier-wave in relation to the amplitude of the message as:


\begin{equation}
    s(t) = A_c\cos{(2 \pi f_c t + 2\pi k_f \int_{0}^{t}(m(\tau)) \,d\tau)}
\end{equation}


Where:

$f_c$ is the carrier frequency

$A_c$ is the carrier amplitude

$m(t)$ is the message signal (which could be anything)


\subsection{Modulation Techniques}

To modulate a FM signal we need a oscilator that can adjust its oscliation frequency based on some voltage input. The 555 Timer IC is usually used for this purpose by putting the 90-degree phase shifted message into the control pin of the timer which will be simulated using arrays and 1D vectors in software

\subsubsection{matricies and vectors in discrete-time digital systems}

computers store numbers one after another, we will use the fameous IEEE 754 double-presicion floats in this paper. In most programming lanugages the indicies start at 0.

\subsubsection{operator overloading clarification}

in C++ and Python its possible to take the cosine of a vector and other simillar operations that translate to doing it for every element in the vector

in other words,

\begin{lstlisting}
    a = np.array([0 pi/2 2*pi])
    b = cos(a)
\end{lstlisting}


turns into:

\begin{lstlisting}
   for i in range(len(a)):
        b[index] = cos(a[index])
\end{lstlisting}


or in C++


\begin{lstlisting}
    std::valarray<double> a {0, pi/2, 2*pi};
    std::valarray<double> b = cos(a);
\end{lstlisting}

turns into:

\begin{lstlisting}
    for(int i = 0; i < a.size(); ++a) \{
        b[i] = cos(a[i]);
    \}
\end{lstlisting}

this is done by the programming language authors and sometimes done manually


\section{Block diagram and software simulation}

\subsection{Overwiew}

\includegraphics[scale=0.15]{block_diagram.png}


block 1 is responsible for modulation and will take the message signal $m(t)$ and output the modulated signal $s(t)$

after modulation during transmission there will be noise which is represented by $w(t)$ and will be simulated using a random variable whose mean is 0 and standard deviation is a small number between 0 and 1 

\subsection{Block1 in Python }

We need to calculate the integral of the signal discretely, to do this we will take a cumulative sum and divide by the sampling frequency to normalize



First we prepare some constants to use later

\begin{lstlisting}
import numpy
fc = 10000
fs = 2*(max(self.fm1, self.fm2, self.fm3) + fc)
time = np.arange(0, 0.1, 1/fs)
fm1 = 20
fm2 = 200
fm3 = 2000
\end{lstlisting}


in this case:


$f_c$ the carrier frequency is set to 10000 Hz,


$f_s$ the sampling frequency is set to double the sum of the highest message tone and the carrier,


the time vector $t$ is created using variable name time between 0 and 1 using our sampling frequency to calculate the increments of time using $\frac{1}{f_s}$

Now we will do the modulation

\begin{lstlisting}
m = cos(2*pi*fm1*time) + cos(2*pi*fm2*time) + cos(2*pi*fm3*time)
m = am * m

s = cos(2*pi*fc*time + 2*pi*np.cumsum(m)/fs)
\end{lstlisting}


%our time vector is multiplied by the scalar $2\pi f_m1$, the cosine of each element is taken 


\subsection{block2}

in block2 we will try to reverse what we did in block1. to do this we will use the hilbert transform and phase derivatives


we will first take the hilbert transform

\begin{equation}
    a(t) = s(t) + js_h(t)
\end{equation}


and take the phase angle

\begin{equation}
    \theta(t) = \arg(a(t)) = \arctan(\frac{s_h(t)}{s(t)})
    = \arctan(\frac{s(t - \frac{pi}{2})}{s(t)})
\end{equation}


fix the arctan using unwrap fucntion in python

\begin{equation}
    \theta_u(t) = unwrap(t)
\end{equation}




now we take the derivative of this to get the instatnaneous phase angle

\begin{equation}
    m(t) = \frac{d\theta_u(t)}{dt}
\end{equation}




\subsection{block2 in python}

this can be done with 1 line of code

NOTE: normalize by dividing with $(2\pi dt)$

NOTE: normalize the dc offset with $-fc$

\begin{lstlisting}
demodulated = np.diff(np.unwrap(np.angle(scipy.signal.hilbert(s))))/(2*pi)*fs - fc
\end{lstlisting}



\subsection{noise}

to prevent noise from causing signficant discomfort a simple low pass filter must be used

\begin{equation}
    LPF(f) = \frac{1}{2 \pi f_c t}
\end{equation}

this $f_c$ means cutoff not carrier


apply this using convolution or frequency domain

\section*{Acknowledgment}

Special thanks to my teachers, friends and the internet and gemini for the help

\section*{References}


\vspace{12pt}
\end{document}
